\documentclass[9pt,letterpaper]{extarticle}
\usepackage[margin=0.7in,top=0.65in,bottom=0.7in]{geometry}
\usepackage[T1]{fontenc}
\usepackage{lmodern}
\usepackage{helvet}
\usepackage{xcolor}
\usepackage{booktabs}
\usepackage{tabularx}
\usepackage{longtable}
\usepackage{colortbl}
\newcommand{\rowsep}{\noalign{\vskip 1.5pt}\arrayrulecolor{black!35}\hline\arrayrulecolor{black}\noalign{\vskip 1.5pt}}
\usepackage{array}
\usepackage{enumitem}
\usepackage[most]{tcolorbox}
\usepackage{microtype}
\usepackage[hidelinks]{hyperref}
\usepackage{needspace}

%% ---------------------------------------------------------------- palette
\definecolor{accent}{HTML}{1F5F8B}   % deep blue
\definecolor{accentlight}{HTML}{E8F0F7}
\definecolor{warn}{HTML}{9A3B1C}     % rust
\definecolor{warnlight}{HTML}{F8ECE6}
\definecolor{keep}{HTML}{2E6B3A}     % green
\definecolor{keeplight}{HTML}{E9F3EB}
\definecolor{ink}{HTML}{222222}
\definecolor{muted}{HTML}{666666}

\color{ink}
\setlength{\parindent}{0pt}
\setlength{\parskip}{4pt}

%% ---------------------------------------------------------------- headings
\makeatletter
\renewcommand\section{\@startsection{section}{1}{\z@}%
  {-10pt \@plus -2pt \@minus -1pt}{3pt}%
  {\sffamily\bfseries\large\color{accent}}}
\renewcommand\subsection{\@startsection{subsection}{2}{\z@}%
  {-8pt \@plus -2pt \@minus -1pt}{2pt}%
  {\sffamily\bfseries\color{accent}}}
\makeatother

%% ---------------------------------------------------------------- boxes
\newtcolorbox{callout}[2][]{enhanced, breakable, colback=accentlight, colframe=accent,
  boxrule=0.6pt, arc=2pt, left=7pt, right=7pt, top=4pt, bottom=4pt,
  fonttitle=\sffamily\bfseries, coltitle=white, colbacktitle=accent,
  title={#2}, #1}
\newtcolorbox{answerbox}[2][]{enhanced, breakable, colback=warnlight, colframe=warn,
  boxrule=0.6pt, arc=2pt, left=7pt, right=7pt, top=4pt, bottom=4pt,
  fonttitle=\sffamily\bfseries, coltitle=white, colbacktitle=warn,
  title={#2}, #1}
\newtcolorbox{keepbox}[2][]{enhanced, breakable, colback=keeplight, colframe=keep,
  boxrule=0.6pt, arc=2pt, left=7pt, right=7pt, top=4pt, bottom=4pt,
  fonttitle=\sffamily\bfseries, coltitle=white, colbacktitle=keep,
  title={#2}, #1}

%% What / Why / How table (longtable breaks across pages)
\setlength{\LTpre}{4pt}
\setlength{\LTpost}{2pt}
\setlength{\arrayrulewidth}{0.3pt}

\newcommand{\hd}[1]{\textbf{#1}}
\let\oldsection\section
\renewcommand{\section}{\needspace{7\baselineskip}\oldsection}
\newcommand{\wwh}{\toprule
\sffamily\bfseries\color{accent} What & \sffamily\bfseries\color{accent} Why & \sffamily\bfseries\color{accent} How \\
\midrule}
\newenvironment{wwhtable}{\footnotesize
\begin{longtable}{@{}>{\raggedright\arraybackslash}p{0.23\linewidth}>{\raggedright\arraybackslash}p{0.355\linewidth}>{\raggedright\arraybackslash}p{0.355\linewidth}@{}}
\wwh}{\bottomrule\end{longtable}}

\begin{document}

%% ---------------------------------------------------------------- title
{\sffamily
{\Large\bfseries\color{accent} What changed, why, and how}\\[3pt]
{\large Checkpoint-2 memo for \emph{Escalate or Answer} (design v3 $\rightarrow$ v4)}\\[2pt]
{\color{muted} For Coen Petto and Cameron Suess \,$\cdot$\, CS 567, Fall 2026 \,$\cdot$\, prepared 29 September 2026}
}

\vspace{4pt}
\noindent{\color{accent}\rule{\linewidth}{1.2pt}}

\begin{callout}{Bottom line}
This is a narrower pass than v2$\to$v3: same population, same corpus
family, same literature, same RQ and MiPP anchor, so nothing here touches
Related Work or the bibliography. What changed is the worked example (now
a real, already-verified Falcon \texttt{-short}/\texttt{-medium} flag
instead of an invented GPU-memory fork), the interaction (a single scored
answer field autofilled after a scripted, paced retrieval trace, accepted
or rejected, instead of a dropdown run card behind a separate source
panel), and the task load (4 forked items per block, no unforked control
items, 12 measured instead of 18). Measures gained a richer per-source,
timed component without giving up the binary H1 test. One tradeoff is
accepted rather than fixed: task set and escalation form are now perfectly
confounded, and that is stated as a limitation, not patched with a Latin
square, for this checkpoint. The proposal is \texttt{new\_proposal\_v4.pdf},
nearer 5.5 pages than v3's five because of a new formal IV/DV/design
statement in the Evaluation Plan; this memo explains each change and what
still needs hand verification.
\end{callout}

\section{Framing and example}
\begin{wwhtable}
Replace the invented GPU-memory fork with a real, already-documented one. &
v3's running example (kestrel-gpu vs.\ peregrine-gpu by model memory) was
plausible but invented for the proposal. The team already had verified
documentation on the Falcon cluster's QoS flag from prior work, and a real,
pre-checked fork is stronger than one built to order. &
The primary worked example is now the Falcon \texttt{-short} vs.\
\texttt{-medium} submission flag, keyed to job wall time, drawn from the
Submitting Jobs page. The abstract, RQ, Section~\ref{sec:fork} analogue and
Figure~1 mock-up all use it; the GPU-partition fork survives only as a
secondary task in the item table. \\
\rowsep
Keep the fork logic, change only the concrete stakes. &
The mechanism (two passages, both right for someone, deciding fact withheld
from the assistant) did not need to change, only which real system it is
grounded in. &
Same definition of a context-dependent fork; same reason the ``newer page
wins'' rule stays dropped; same reason the agent cannot ask up front. \\
\end{wwhtable}

\section{Interaction model}
\begin{wwhtable}
Replace the dropdown run card and separate source panel with a scripted
chat trace and one scored answer field. &
Four dropdowns (one scored, the rest uncontested decoys) behind a card was
harder to build and to explain to a participant as ``one decision.'' A
single field that the assistant fills, which the participant then accepts
or rejects, makes the accept/reject act itself the primary logged behavior
instead of a submitted card. &
The assistant prints a persistent, scripted thinking-and-retrieval trace,
paced 3--5\,s per line, that names \emph{Source A} and \emph{Source B} as
clickable chips; a scripted decision line fills the one answer field; the
participant clicks \textbf{Accept} or \textbf{Reject} (Reject opens a
two-option picker to pick the alternate, which doubles as one-step undo). \\
\rowsep
Fold the source panel into the trace; no separate panel. &
A right-hand source panel duplicated what the trace already cited, and
having two places to look at sources made it harder to say precisely where
verification happens. &
Clicking a chip opens that passage in an overlay over the trace with the
relevant lines highlighted; closing it returns to the trace. The trace is
now the only place a source can be opened, so ``verified'' is unambiguous:
opened a source from the trace between Send and Accept. \\
\rowsep
Rename the forms to match the new mechanic. &
``Run card of dropdowns'' and ``the card is pre-filled'' language no longer
described anything; the forms still need the same announced/blocking
definition. &
S1 modal stop (trace freezes, participant must pick), S2 inline warning
(trace continues, assistant decides, glyph plus recommendation shown), S3
provenance-only (trace continues, assistant decides, citation only). Same
announced/blocking logic as v3's modal/inline/provenance-only, renamed
S1--S3 and drawn for the Falcon example in the new Figure~1. \\
\end{wwhtable}

\section{Task load and controls}
\begin{wwhtable}
Drop the unforked control items. &
v3 used 2 no-fork control items per block (of 6) for a baseline
page-opening rate and so two citation chips never signalled a fork by
themselves. Every v4 trace already cites two sources on every task, so a
citation count can no longer leak the fork either way, and the
provenance-only form is itself the baseline for how often people check
when nothing announces anything. &
4 items per block, all forked, in every block: no control items. This is
disclosed to participants up front (the instructions state the assistant
can be wrong on any task), so the fork can no longer be ``found'' by
contrast with unflagged items; that contrast is gone as a source of
information for participants, and as a design lever for the researchers. \\
\rowsep
Shrink the item count. &
Fewer, more heavily instrumented items are enough for the same number of
measured forks once controls are cut, and a shorter block shortens the
session. &
12 measured tasks (4/block $\times$ 3 blocks, 6 with a correct
recommendation and 6 with a wrong one) plus 2 practice tasks, down from 18
measured (12 flagged $+$ 6 control) plus 2 practice in v3. \\
\end{wwhtable}

\section{Measures}
\begin{wwhtable}
Keep the binary verification measure for H1; add a richer per-source, timed
measure alongside it. &
v3's ``verified'' (opened a diverging page, yes/no, anchored to the item)
stays exactly comparable across v3 and v4 and is still what H1 is tested
on. But it collapses ``checked nothing,'' ``checked only the recommended
source,'' and ``checked both'' into the same value, which loses the more
diagnostic question of \emph{what} people check, not just whether. &
\emph{Verified} (binary, unchanged in spirit) remains the primary DV for
H1. New, per task: \emph{Source pattern} (recommended only, non-recommended
only, both, or neither opened) and \emph{Dwell} (seconds each source was
open, summed over openings), reported descriptively by form and analyzed
alongside \emph{Followed} for appropriate reliance. \\
\rowsep
State clearly that this is additive. &
A reviewer could otherwise read the new measures as a quiet replacement of
the original hypothesis test. &
The proposal keeps H1 phrased on Verified exactly as before; Source pattern
and Dwell are named secondary and descriptive, not substituted into the
formal H1/H2/H3 tests. \\
\end{wwhtable}

\section{Spec sheet, not a per-item situation line}
\begin{wwhtable}
Externalize the deciding fact into a persistent, block-level spec sheet. &
v3's deciding fact lived in a one- or two-line ``Your situation'' field on
each task card, shown once per item. That does not match the cover story
(a student completing an assignment on Falcon) and buries the mechanism
inside the item rather than making it a standing, checkable object. &
Each block is now framed as one assignment; a spec sheet of plain
constraint statements (``Stage 2, training: about 30\,h, 1 GPU, 40\,GB GPU
memory'') stays on screen for the whole block. The assistant never
receives it; the participant does. \\
\rowsep
Keep the sheet easy to check and easy to skip. &
The sheet is what turns ``did they verify'' into a measure of complacency
rather than an artifact of a fact being forced in front of them on every
single item, the way the per-item situation line did. &
Plain constraint statements only, no scoring hints; the same sheet backs
multiple tasks in a block, so a participant who never reads it is making
the same mistake repeatedly rather than being re-prompted each time. \\
\end{wwhtable}

\section{Known limitation: task set is confounded with escalation form}
\begin{wwhtable}
Accept, and state plainly, that each task set is now tied to one system. &
v3 crossed task set and form in a $3\times3$ Latin square, so any
difference between forms could not be attributed to set content. v4 ties
S1 always to the job-submission set, S2 always to access/connection, S3
always to storage/software. Task-set difficulty is now perfectly
collinear with escalation form: any effect between forms could be driven
by which set of tasks a form happened to get, not by the form itself. &
Name it in Threats as a stated, deliberate limitation for this
checkpoint's design, not an oversight: crossing set with form (as in v3's
Latin square) would remove the confound at no measurement cost, and the
team considered it, but decided against restructuring for this
checkpoint. The matching procedure (sets matched on trace length, fork
line index, query--source overlap band, fork type) and the task random
effect in every model are the only defenses that remain, and the memo does
not propose the Latin-square fix as something to do now. \\
\rowsep
Say what it costs. &
A reviewer can no longer be told that an observed form effect is
attributable to form alone. &
The proposal's own Threats paragraph says so directly; this memo does not
soften that. \\
\end{wwhtable}

\section{Length, formal design statement, and citations}
\begin{wwhtable}
Add a MacKenzie-style formal IV/DV/design-statement block to the
Evaluation Plan before the narrative procedure. &
v3's Evaluation Plan stated the design in prose only. A reviewer expects
independent variables, dependent variables and a formal design statement
written out explicitly, which is more rigorous but longer. &
Design section now opens with Escalation Form and Recommendation
Correctness spelled out as IVs with their levels, Block Order as a
between-subjects nuisance variable, a one-sentence formal design statement
($3\times2$ within-subjects, two replicates per cell, 12 tasks/participant,
216 tasks / 36 per cell at $N{=}18$), and DVs (a)--(g) listed individually
before the narrative procedure resumes. \\
\rowsep
Accept the length tradeoff; flag where to cut if the budget is hard. &
The extra rigor pushed the proposal from about 5 pages (v3) to nearer 5.5
(v4). Related Work was left untouched in this pass and is the one section
with slack, since nothing about the literature or its citations changed. &
If page count is a hard constraint, trim Related Work first; do not cut
the new IV/DV/design-statement block, since that is exactly what a formal
reviewer will look for. \\
\rowsep
Bibliography and citations: unchanged. &
No new claims, sources or citation audit were needed for this pass. &
\texttt{references.bib} and every in-text citation are identical to v3;
see \texttt{what\_changes.tex} for that audit rather than repeating it
here. \\
\end{wwhtable}

\begin{keepbox}{What v4 keeps from v3}
\begin{itemize}[leftmargin=1.2em, itemsep=1pt]
  \item \hd{The RQ, title and MiPP anchor}, unchanged word for word.
  \item \hd{The CSU student population and the Falcon/SNA corpus family},
  same recruitment route and screening logic.
  \item \hd{Within-subjects design with all six block orders}, NASA-TLX,
  the intrusion item, the manipulation check and the end-of-session
  ranking with a free-text reason.
  \item \hd{The context-dependent fork definition}, the FictionalQA
  blind-vs-informed filter, and the reason mid-task escalation is used
  instead of asking up front.
  \item \hd{The announced/blocking definition of the three forms} and the
  MiPP C1--C3 constraints held constant across them.
  \item \hd{The agentic drafting loop and its fork check}, frozen and
  replayed identically to every participant.
  \item \hd{Related Work and the bibliography}, untouched in this pass.
\end{itemize}
\end{keepbox}

\section{Prioritized checklist for checkpoint 2}
\begin{enumerate}[leftmargin=1.6em, itemsep=2pt, label=\sffamily\bfseries\color{accent}\arabic*.]
  \item \hd{Verify the \texttt{-short}/\texttt{-medium} flag syntax and
  wall-time limits} against the frozen Submitting Jobs page; the current
  text uses \texttt{-short}/\texttt{-medium} as written on the whiteboard
  and has not been checked against \texttt{--qos=} syntax or the actual
  limits.
  \item \hd{Draft the second source pair for the S2 and S3 sets} (tasks
  2.3--2.4 and 3.3--3.4) from the drafting pipeline once the snapshot is
  frozen; do not hand-seed them.
  \item \hd{Decide, and write down the decision,} on the task-set/form
  confound: confirm the sets are tied to systems as agreed on the
  whiteboard, and if the tie is kept, make sure the matching procedure and
  the task random effect are documented in Threats as the stated defenses.
  \item \hd{Read \texttt{new\_proposal\_v4.tex} end to end} against this
  memo and change anything that is not your view, in particular: the
  Falcon example, the 12-task/no-control design, the spec sheet wording,
  and the new IV/DV/design-statement block.
  \item \hd{If the page budget is hard, trim Related Work first}; do not
  cut the new formal design-statement block to save space.
  \item \hd{Confirm the Team paragraph roles} and the Use-of-AI note about
  this revision pass before submitting.
\end{enumerate}

\end{document}
